\subsection*{Sequentialism and Parallelism}

Repetitive tasks are common in programming.

Currently, the most widely used parallelization paradigm is \textbf{Single Program, Multiple Data (SPMD)}. In this paradigm, code is written once but executed as multiple instances simultaneously. Each instance operates independently unless synchronization is explicitly implemented. Croqtile adopts this paradigm and abstracts parallelism using the \textbf{parallel-by} block.

A sequential loop is needed when later work depends on earlier work, such as accumulating successive K slices. Croqtile separates the bounded range declared by \co@with-in@ from the traversal performed by \co@foreach@. This also permits work before and after the traversal, as used in buffered pipelines.

The following sections explain how Croqtile constructs loops and expresses parallel work.
